% ===================================================================
\section{Mixed reasoning drill}
% ===================================================================
Statement- and assertion--reason-style items across the course. For assertion--reason items use the usual key:
\textbf{(A)} both true and R explains A; \textbf{(B)} both true, R does not explain A; \textbf{(C)} A true, R false;
\textbf{(D)} A false, R true; \textbf{(E)} both false. Answers with reasons follow.

\subsection*{Questions}
\begin{enumerate}[label=\textbf{Q\arabic*.},leftmargin=2.6em]
\item \textbf{A:} Synchronous calving lowers the fraction of wildebeest calves lost to predators. \textbf{R:} Predators have a
limited daily capacity to kill and eat.
\item Which statements about a Leslie matrix population with constant vital rates are true? (i) The stable age distribution
is independent of the initial vector. (ii) At the SAD every class stays the same size. (iii) If only one age class
reproduces, the population may never settle to a SAD. (iv) $\lambda>1$ exactly when $R_0>1$.\\
(a) i, iii, iv\quad (b) i, ii\quad (c) ii, iv\quad (d) all four
\item \textbf{A:} Protecting loggerhead eggs and hatchlings cannot reverse the population decline. \textbf{R:} The sensitivity of
$\lambda$ is largest for the egg stage.
\item In the Ricker model with $r=1.5$ and $K=1000$, the population will\\
(a) approach $K$ monotonically\quad (b) approach $K$ with damped oscillations\quad (c) settle on a 2-cycle\quad (d) behave
chaotically because $K$ is large
\item \textbf{A:} The continuous logistic model cannot produce population cycles. \textbf{R:} Its per-capita growth rate falls
linearly with density.
\item Hassell et al.\ (1976) found 22 of 24 insect populations in the monotone-damping region. The best reading is:\\
(a) chaos is impossible in insects\quad (b) insects have low growth rates\quad (c) a one-lag model with a rigid form may be
biased toward stability\quad (d) integer counts destroyed the chaos in those data
\item \textbf{A:} Integerised population models cannot be chaotic. \textbf{R:} A deterministic map on a finite set of states is
eventually periodic.
\item Two competitors: $K_1=100$, $K_2=80$, $\alpha_{12}=0.5$, $\alpha_{21}=1.5$. Then\\
(a) stable coexistence because $\alpha_{12}\alpha_{21}<1$\quad (b) species 1 wins\quad (c) species 2 wins\quad (d) priority
effect
\item \textbf{A:} In Park's experiments the species with the larger single-species equilibrium sometimes lost. \textbf{R:} The
outcome of competition depends on competition coefficients as well as carrying capacities.
\item Which is \emph{not} evidence that niche partitioning allows coexistence?\\
(a) \emph{P.\ caudatum} and \emph{P.\ bursaria} feed at different depths\quad (b) the golden spiny mouse becomes nocturnal
when the common spiny mouse is removed\quad (c) \emph{P.\ aurelia} excludes \emph{P.\ caudatum}\quad (d) finch beaks differ
more where species co-occur
\item In the basic L--V predator--prey model, doubling the prey's intrinsic growth rate $r$\\
(a) doubles $H^*$\quad (b) doubles $P^*$\quad (c) doubles both\quad (d) changes neither
\item \textbf{A:} Gause's \emph{Paramecium}--\emph{Didinium} cultures oscillated only with regular immigration. \textbf{R:} Left to
itself the simple system went extinct.
\item \textbf{A:} Huffaker obtained sustained predator--prey oscillations only with a complex arrangement of oranges.
\textbf{R:} Prey dispersing faster than predators could colonise patches ahead of them.
\item A Batesian mimic becomes much commoner than its model. Expected result:\\
(a) both are better protected\quad (b) predators learn to attack the signal more often\quad (c) the model becomes a mimic\quad
(d) nothing changes, because the signal is unchanged
\item \textbf{A:} Restoring a shallow lake usually needs a lower phosphorus level than the one at which it collapsed.
\textbf{R:} The vegetation--turbidity feedback loop is reinforcing.
\item Which restoration method is unsuitable for a shallow, unstratified lake?\\
(a) sediment dredging\quad (b) biomanipulation\quad (c) hypolimnetic oxygenation\quad (d) phosphorus inactivation
\item \textbf{A:} Lake Washington's transparency rose to about 10\,m in 1976. \textbf{R:} Phosphorus loading was diverted from the
lake in 1976.
\item A population has $b=1.25$, $d=1.0$ per capita. Starting from 4 individuals, its probability of eventual extinction is
about\\
(a) 0\quad (b) 0.41\quad (c) 0.8\quad (d) 0.2
\item \textbf{A:} The passenger pigeon kept declining after commercial hunting stopped. \textbf{R:} Social facilitation in
finding food made per-capita success fall at low numbers.
\item Which pair is correctly matched?\\
(a) Muir -- resource conservation\quad (b) Pinchot -- greatest good for the greatest number for the longest time\quad
(c) Leopold -- preservation for spiritual communion\quad (d) Brandis -- land ethic
\item \textbf{A:} Cuvier's work removed the belief that no species ever went extinct. \textbf{R:} Cuvier accepted that species
change into new species.
\item \textbf{A:} Under blending inheritance natural selection would soon have nothing to act on. \textbf{R:} Blending halves the
heritable variance every generation.
\item A sample at one locus fits HW proportions ($\chi^2=0.8$). The best conclusion is:\\
(a) the locus is neutral\quad (b) mating is random in the population\quad (c) there is no detectable departure from HWE in
this sample\quad (d) the population is large
\item A recessive condition affects 1 in 2500. The proportion of carriers among unaffected people is closest to\\
(a) 0.02\quad (b) 0.039\quad (c) 0.08\quad (d) 0.0004
\item \textbf{A:} Natural selection acts on individuals but populations evolve. \textbf{R:} An individual's genotype does not change
during its life.
\item \textbf{A:} After the 1977 drought, chicks on Daphne Major had deeper beaks than before. \textbf{R:} Beak depth is
heritable ($h^2\approx0.65$) and deep-beaked birds survived better.
\item Which statement about selection is correct?\\
(a) it produces variation when needed\quad (b) it tends toward perfection\quad (c) it can only act on existing variation, though
mutation and recombination supply new variants\quad (d) it requires that some individuals die
\item \textbf{A:} A rare, advantageous recessive allele spreads very slowly at first. \textbf{R:} When rare, most of its copies
are in heterozygotes, where it is not expressed.
\item \textbf{A:} The Extended Evolutionary Synthesis is a refutation of the Modern Synthesis. \textbf{R:} Some acquired
epigenetic states can be inherited across generations.
\item Which pair of cases from the course both show a threshold beyond which reversing the original change is not enough?\\
(a) Lake Veluwe and the passenger pigeon\quad (b) Gause's exclusion and the hare--lynx cycle\quad (c) the Leslie SAD and
the logistic\quad (d) TEDs and Project Tiger's early increase
\end{enumerate}

\subsection*{Answers and reasons}
\begin{enumerate}[label=\textbf{Q\arabic*.},leftmargin=2.6em]
\item \textbf{A.} A capped daily kill means losses $\le Kw$; packing births into a short window lowers the fraction.
\item \textbf{(a).} (ii) is false: proportions are constant, numbers grow by $\lambda$. (iii) is the imprimitive case.
\item \textbf{C.} A is true; R is false: egg-stage elasticity is low; large juveniles have the highest elasticity.
\item \textbf{(b).} $g'(K)=1-r=-0.5$: negative with modulus $<1$, so damped oscillations; $K$ does not matter.
\item \textbf{B.} Both true, but the reason a 1-D ODE cannot cycle is that it moves along a line and cannot reverse at a point;
the linear form of the rate is not the explanation (any smooth one-variable ODE cannot oscillate).
\item \textbf{(c).} Turchin and Taylor's point. (a) overstates; (b) contradicts insect fecundity; (d) is a different argument
about models, not those fits.
\item \textbf{A.} Pigeonhole argument (slide 33).
\item \textbf{(b).} Can 2 invade 1? $\alpha_{21}K_1=150>K_2=80$: no. Can 1 invade 2? $\alpha_{12}K_2=40<K_1=100$: yes. So
species 1 wins (its isocline, intercepts 100 and 200, lies outside species 2's, 53.3 and 80). The trap is (a):
$\alpha_{12}\alpha_{21}=0.75<1$ is not sufficient.
\item \textbf{A.}
\item \textbf{(c).} Exclusion on one niche is Gause's principle, not partitioning.
\item \textbf{(b).} $P^*=r/p$; $H^*=m/(ap)$ is unchanged.
\item \textbf{A.}
\item \textbf{A.} Differential dispersal through a patchy landscape gives metapopulation persistence.
\item \textbf{(b).} The signal loses its information when most bearers are palatable.
\item \textbf{A.} The reinforcing loop with saturating ends gives two stable states and hence hysteresis.
\item \textbf{(c).} No hypolimnion to oxygenate.
\item \textbf{C.} Diversion was completed in 1968; the 1976 jump followed the return of \emph{Daphnia}.
\item \textbf{(b).} $(1/1.25)^4=0.8^4=0.41$.
\item \textbf{A}, with a caveat: the slide calls habitat loss plus social facilitation the leading hypothesis (``we do
not know''), so R is the course's best explanation rather than an established fact.
\item \textbf{(b).}
\item \textbf{C.} Cuvier rejected change of species.
\item \textbf{A.} Jenkin's argument.
\item \textbf{(c).} Non-rejection does not prove any assumption.
\item \textbf{(b).} $q=1/50=0.02$; $2pq/(1-q^2)=2(0.98)(0.02)/0.9996=0.0392$.
\item \textbf{A.}
\item \textbf{A.} All four postulates satisfied.
\item \textbf{(c).} (d) is false: differences in reproduction suffice.
\item \textbf{A.}
\item \textbf{D.} A is false (it is an extension, nothing contradicts the MS); R is true (vinclozolin, toadflax).
\item \textbf{(a).} Hysteresis in the lake; a population below its MVP not recovering after hunting stopped.
\end{enumerate}

\subsection*{Long-answer prompts (20-mark style)}
\begin{enumerate}[leftmargin=2.2em]
\item Why is biology hard to formalise, and how does BI2113 respond? Use the wildebeest to illustrate all four features.
\item ``Where should conservation effort go?'' Answer with the loggerhead matrix, elasticity and TEDs; then contrast with the
tiger (reserve design, counting) and the passenger pigeon (thresholds).
\item Simple models predict chaos at high growth rates, yet it is almost never seen. Explain with both resolutions.
\item Competition theory predicts exclusion in most cases, yet similar species coexist. Resolve with the model and experiments.
\item Real predator--prey cycles exist, but the L--V explanation is not accepted. Discuss with Gause, Huffaker and Krebs.
\item Why does phosphorus reduction often fail to restore shallow lakes? Include the feedback loop, hysteresis, Lake
Washington and the restoration toolkit.
\item Trace evolutionary thought from fixity to the Modern Synthesis, and assess whether an Extended Synthesis is needed.
\item Derive Hardy--Weinberg with its assumptions; explain why it mattered for Darwin and what each departure does.
\item ``Natural selection is not evolution.'' Discuss with the Daphne Major finches.
\item Thresholds and hysteresis appear in lakes, small populations, competition and genetics. Compare four cases.
\end{enumerate}
